\section{Bounding the Smallest Leontovich Graph}
\label{sec:bounds}

\subsection{Exhaustive sweep}

We tested all 1,018,690,329 connected simple graphs on $m \le 11$
vertices~\cite{oeis_A001349} (generated by McKay's
\texttt{geng}~\cite{mckay2014practical}) against the
$E_n^{(d)}$ filter for $n \le 200$, $d \le 20$.

\begin{compresult}[Computational sweep certificate]
  No connected simple graph on $m \le 11$ vertices produced a near-path
  violation $\Hom(E_n^{(d)},H)<\Hom(P_n,H)$ in the tested range
  $n \le 200$, $d \le 20$. For $m \le 9$, this was supplemented by
  brute-force checks over all source trees on $n \le 15$ vertices.
\end{compresult}

This computationally bounds the tested near-path obstruction in
Problem~4.3 of Galvin et al.~\cite{galvin2025}; it should not be read as an
unconditional proof that no other source-tree shape can beat the path for
$m=10,11$. Because our search space comprised simple graphs (generated by
\texttt{geng}), this evidence also does not exclude the possibility of a
looped or general graph on fewer than 12 vertices being strongly Leontovich
in the sense of~\cite{galvin2026}.
Galvin et al.~\cite{galvin2025} established an upper bound of
$|V| = 961{,}601$ using 4-orbit symmetric trees with a single looped
root. Expanding the degrees of freedom to 5-orbit looped symmetric trees
gives a smaller verified strongly Leontovich example on $1{,}822$ vertices
at $\hat{T}(1, 35, 1, 50)$, which will be used below because its stronger
margin is convenient for the double-cover construction. Furthermore, since our exhaustive search
includes all simple non-bipartite graphs (those containing odd cycles), this
sweep gives lower-bound evidence for simple strongly Leontovich graphs within
the tested near-path regime, significantly narrowing the uncertainty gap from
the bottom up. It cannot be upgraded directly to an unconditional small-target
nonexistence theorem: Section~\ref{sec:doublecover} exhibits a first crossover
only at $n = 17{,}340$, far beyond the sweep horizon used here.

\begin{table}[H]
\centering
\caption{Exhaustive sweep: all connected $H$ with $|V(H)| \le 11$.}
\label{tab:exhaustive}
\begin{tabular}{@{}ccccc@{}}
  \toprule
  $m$ & Connected $H$ & Cumulative & Operations & Violations \\
  \midrule
  1 & 1       & 1       & ---                     & \textbf{0} \\
  2 & 1       & 2       & ---                     & \textbf{0} \\
  3 & 2       & 4       & ---                     & \textbf{0} \\
  4 & 6       & 10      & $1.8 \times 10^7$   & \textbf{0} \\
  5 & 21      & 31      & $9.9 \times 10^7$   & \textbf{0} \\
  6 & 112     & 143     & $7.6 \times 10^8$   & \textbf{0} \\
  7 & 853     & 996     & $7.9 \times 10^9$   & \textbf{0} \\
  8 & 11,117  & 12,113  & $1.3 \times 10^{11}$ & \textbf{0} \\
  9 & 261,080 & 273,193 & $4.0 \times 10^{12}$ & \textbf{0} \\
  10 & 11,716,571 & 11,989,764 & 45\,s$^\dagger$ & \textbf{0} \\
  11 & 1,006,700,565 & 1,018,690,329 & 136\,min$^\dagger$ & \textbf{0} \\
  \bottomrule
\end{tabular}
\smallskip

\noindent{\footnotesize $^\dagger$For $m \le 9$, every tree on
$n \in \{5,\ldots,15\}$ vertices was tested (brute-force); the
Operations column counts hom evaluations. For $m = 10, 11$, only
near-path trees $E_n^{(d)}$ ($n \le 200$, $d \le 20$) were used
via the $E_n^{(d)}$ filter (wall-clock time shown). These rows are
bounded filter results in the sense of Remark~\ref{rem:filter}, not
full all-tree sweeps.}
\end{table}

\subsection{Asymptotic defense via Perron--Frobenius}

\begin{remark}[Filter heuristic]
  \label{rem:filter}
  The near-path family $E_n^{(d)}$ is a natural high-sensitivity filter for
  path-minimizer failures.
  The path $P_n$ achieves the minimal exponential growth rate
  $\lambda_1(H)^n$ (Csikv\'ari and Lin~\cite{csikvari2015}). For any
  connected, non-regular graph $H$, the Perron--Frobenius eigenvector equation guarantees
  $\Delta(H) > \lambda_1(H)$ strictly. Trees with macroscopic branching
  (such as stars) grow at an exponential rate strictly greater than $\lambda_1(H)^n$,
  exponentially faster than the path.
  To challenge the path, a tree must match the growth base $\lambda_1^n$
  while achieving a smaller leading coefficient; near-path trees
  $E_n^{(d)}$ provide a controlled way to test such localized boundary
  defects. We use this as a computational filter, not as an unconditional
  proof that all possible source-tree obstructions have been exhausted.
\end{remark}

This reduces the per-target cost from brute-force $O(t \cdot n \cdot m^2)$
(over $t$ source trees) to $O(d_{\max} \cdot n_{\max} \cdot m)$ via
$O(m)$ matrix--vector iteration up to $n = 200$, $d \le 20$.
Filtering across all 1,018,690,329 connected graphs with $m \le 11$
yields zero near-path violations in the tested range.
The exact exhaustive and certified-table computations in this paper use
integer arithmetic, so compiler flags such as \texttt{-march=native}
change runtime but not the certified values themselves; only
floating-point screens are treated as conditional filters before exact
verification.

\begin{table}[H]
\centering
\caption{Conditional results from the $E_n^{(d)}$ filter and bipartite sweeps.}
\label{tab:conditional}
\begin{tabular}{@{}lll@{}}
  \toprule
  Domain & Scope & Outcome \\
  \midrule
  All connected, $m \le 11$ & $n \le 200$, $d \le 20$ & no near-path hit \\
  Bipartite, $m \le 14$ & $n \le 199$, $d \le 20$ & no near-path hit \\
  Bipartite, $m = 15$ & $n \le 199$, $d \le 20$ & $H^*$ exists \\
  Bipartite depth-2, $m \le 17$ & $d=2$ sweep & no hit below $H_{18}$ \\
  Trees, $m \le 23$ & tested family & no hit \\
  Symmetric trees, $|V| \le 77$ & tested family & no hit \\
  \bottomrule
\end{tabular}
\end{table}

% ════════════════════════════════════════════════════════════════════════
